% Abstract (250 words max for Nature Machine Intelligence)

Parkinson's disease (PD) exhibits profound heterogeneity in progression patterns, clinical presentation, and treatment response, necessitating precision medicine approaches. However, integrating diverse data modalities—clinical assessments, neuroimaging biomarkers, and genetic profiles—into interpretable predictive models remains challenging. We present the Graph-Informed Multimodal Attention Network (\giman), a comprehensive framework that combines graph neural networks with explainable AI to address three critical tasks: (1) identifying longitudinal progression subtypes, (2) predicting prodromal-to-clinical conversion, and (3) providing clinically interpretable explanations. 

Using 536 PD patients and 194 prodromal subjects from the Parkinson's Progression Markers Initiative, \giman\ constructs patient similarity graphs based on multimodal features and employs Graph Attention Networks (GATs) with multi-head attention to learn disease representations. For longitudinal analysis (Phase 4), trajectory embedding via Variational Autoencoders followed by clustering revealed three distinct progression subtypes: fast motor decline (22\%, 7.8 UPDRS points/year), moderate (54\%, 3.2 points/year), and cognitive-first (24\%, -1.8 MoCA points/year). These subtypes showed distinct biomarker profiles (DaTscan SBR, CSF $\alpha$-synuclein, GBA mutations) and enabled 38\% reduction in clinical trial sample sizes. For prodromal prediction (Phase 5), Cox and DeepSurv models achieved 0.79 C-index, identifying baseline UPDRS-III, sex, and RBD as key predictors.

Critically, we developed a six-method explainability framework (Phase 6) combining attention visualization, GNNExplainer, gradient-based attribution (IntegratedGradients, GradientSHAP), embedding clustering, and counterfactual analysis. Cross-method validation demonstrated 88-95\% consensus on feature importance, with attention weights revealing clinically meaningful patient similarity. \giman\ provides an end-to-end, interpretable framework for PD precision medicine, directly applicable to clinical trial design and personalized treatment strategies.
